\documentclass[../../main.tex]{subfiles}

\begin{document}


\section{Entanglement entropies}
\label{sec:entanglement-entropies}

Entanglement can be quantified in various different fashions. Here, we want to explore the von Neumann and the Rényi entropies.

\subsection{The von Neumann entanglement entropy}
The von Neumann entropy is defined as
\begin{equation}
      S(\rho_A)\equiv-\Tr\bigl(\rho_A\ln\rho_A\bigr)
\end{equation}
$S$ is zero for two unentangled systems and non-zero for entangled systems.
\begin{claim}
      $S(\rho_A)\equiv-\Tr\bigl(\rho_A\ln\rho_A\bigr) = -\sum_i\lambda_i\ln\lambda_i =-\sum_i s_i^2\ln s_i^2.$
\end{claim}
\begin{claimproof}
      Since $\rho_A$ is hermitian, we can use the spectral theorem to write it as
      $\rho_A=\sum_i\lambda_i P_i$ where $\lambda_i$ are its eigenvalues and the eigenvectors
      $\ket{u_i}$ give $P_i=\dyad{u_i}{u_i}$, since
      \begin{equation}
      \rho_A=UDU^\dagger=\bigl(\ket{u_1},\dots,\ket{u_n}\bigr)
      \begin{pmatrix}\lambda_1& &\\ &\ddots&\\ & &\lambda_n\end{pmatrix}
      \begin{pmatrix}\bra{u_1}\\\vdots\\\bra{u_n}\end{pmatrix},
      \qquad\braket{u_i}{u_j}=\delta_{ij}.
      \end{equation}
      For a function $f$ defined on the spectrum of $\rho_A$, the spectral
      functional calculus\footnote{See for example Nicholas, J. H. (2008). \emph{Functions of matrices: theory and computation.}} gives
      \begin{equation}
      f(\rho_A)\equiv\sum_i f(\lambda_i)\dyad{u_i}{u_i}
      =U\diag\bigl(f(\lambda_i)\bigr)U^\dagger.
      \end{equation}
      Therefore, one can rewrite the von Neumann entanglement entropy as
      \begin{align}
      S(\rho_A)&\equiv-\Tr\bigl(\rho_A\ln\rho_A\bigr) \\
      &=-\Tr\bigl(\rho_A\,U\ln(\diag(\lambda_i))\,U^\dagger\bigr)\notag\\
      &=-\Tr\bigl(U\diag(\lambda_i)\underbrace{U^\dagger U}_{=\mathds 1}\ln(\diag(\lambda_i))U^\dagger\bigr)\notag\\
      &=-\sum_i\lambda_i\ln\lambda_i
      =-\sum_i s_i^2\ln s_i^2.
      \end{align}
\end{claimproof}



\paragraph{Example}

In our case, for example, $L=2$, $N_f=1$, $N_\mathrm{ph}=1$:
\begin{align}
& \{\ket{0},\ket{1}\}\ \&\ \{\ket{01},\ket{10}\}
\quad\text{as bases of }\mathcal H_b\ \&\ \mathcal H_f,\\
&\Rightarrow\{\ket{0, 01};\ \ket{0, 10};\ \ket{1, 01};\ \ket{1, 10}\}
\quad\text{as basis of }\mathcal H_b\otimes\mathcal H_f,\\
&\Rightarrow\ket{\psi}=\sum_{i=1}^{2}\sum_{j=1}^{2}\psi_{ij}\ket{i}_b\ket{j}_f
\quad\text{with }(\psi)_{ij}=\psi_{ij}\ \&\ \psi\in\mathbb C^{2\times2},\\
&\Rightarrow S(\rho_A)=-\sum_\ell s_\ell^2\ln s_\ell^2.
\end{align}
Note that since $\Trace\rho^{AB}=\Trace\rho^A=\Trace\rho^B=1$ we must have
$\sum_\ell\lambda_\ell=\sum_\ell s_\ell^2=1$.

\subsection{The Rényi entanglement entropies}

The Rényi entropy of order $n$ is defined as
\begin{align}
      S_n(\rho_A) &\equiv \frac{1}{1-n} \ln (\Trace[\rho_A^n]) \\
      &=  \frac{1}{1-n} \ln \left( \sum_{i=1}^{d_A} \lambda_i^n \right) \equiv S_n(\vec \lambda),
\end{align}
using again, that any hermitian matrix is diagonalizable and the cylic property of the trace. Keep in mind, that $0 \leq \lambda_i \leq 1$ and $\sum_i \lambda_i = 1$.

The Rényi entropy has the following useful properties, which we will prove subsequently:
\begin{itemize}
      \item The Rényi entropy of order one reduces to the von Neumann entropy.
      \item $\lim_{n \to \infty} S_n (\vec \lambda) \equiv S_\infty (\vec \lambda)= - \ln \max_i \lambda_i$ and $S_2(\vec \lambda) \leq 2 S_\infty(\vec \lambda)$.
      \item We have $S_n \geq S_m$ for $n<m$ and therefore $S_2$ is a lower bound for the Von Neumann entropy $S_1$.
\end{itemize}

\begin{claim}
The Rényi entropy of order one reduces to the von Neumann entropy.
\end{claim}
\begin{claimproof}
      We can examine the first order by using l'Hôpital's rule, which yields
      \begin{align*}
            S_1(\rho_A) &= \lim_{n \to 1} \frac{\ln (\Trace[\rho_A^n])}{1-n} 
            \overset{\mathrm{H}}{=} \lim_{n \to 1} \frac{\left( \ln \sum_{i=1}^{d_A} \lambda_i^n \right)'}{(1-n)'} \\
            &= - \lim_{n \to 1} \frac{\sum_{i=1}^{d_A} \lambda_i^n \ln \lambda_i }{\sum_{i=1}^{d_A} \lambda_i^n} 
            = - \frac{\sum_{i=1}^{d_A} \lambda_i \ln \lambda_i }{\sum_{i=1}^{d_A} \lambda_i} \\
            &= - \sum_{i=1}^{d_A} \lambda_i \ln \lambda_i .
      \end{align*}
\end{claimproof}

\begin{claim}
      $\lim_{n \to \infty} S_n (\vec \lambda) \equiv S_\infty (\vec \lambda)= - \ln \max_i \lambda_i$ and $S_2(\vec \lambda) \leq 2 S_\infty(\vec \lambda)$.
\end{claim}
\begin{claimproof}
      Order the different eigenvalues of $\rho_A$ such that $\lambda_1 \leq \lambda_2 \leq \ldots \leq \lambda_2$. Furthermore, denote the number of degenerate $\lambda_1\equiv\lambda_{\text{max}}$ as $m$. It then follows      
      \begin{align*}
            S_\infty (\vec \lambda) &= \lim_{n \to \infty} \frac{1}{1-n} \ln \left( \sum_{i=1}^{d_A} \lambda_i^n \right) \\
            &= \lim_{n \to \infty} \frac{1}{1-n} 
            \ln \left(  
            \sum_{i=1}^{m} \lambda_\text{max} ^n + \sum_{i=m+1}^{d_A}  \lambda_i^n
            \right) \\
            &= \lim_{n \to \infty} \frac{1}{1-n} 
            \ln \left(  m \lambda_\text{max}^n + 0 \right) \\
            &= \lim_{n \to \infty} \frac{\ln m + n \ln \lambda_\text{max}}{1-n} = - \ln \max_{i} \lambda_i .
      \end{align*}
      The second inequality quickly follows
      \begin{equation*}
            S_2 (\vec \lambda) = - \ln(\sum_{i=1}^{d_A} \lambda_i^n) \leq - \ln (\lambda_\text{max}^2) = -2 S_\infty (\vec \lambda),
      \end{equation*}
      since the logarithm decreases monotonically.
\end{claimproof}

\begin{claim}
      $S_n \geq S_m$ for $n<m$.
\end{claim}
\begin{claimproof}
      Define $z_i(n) = \lambda_i^n / \sum_{j=1}^{d_A} \lambda_j^n$, which is a probability distribution, and $F(n) = \ln \sum_{j=1}^{d_A} \lambda^n_j$. We take n to be a continuous variable and differentiate $F$ with respect to it using again that $\frac{d}{dn} \lambda^n = \lambda^n \ln \lambda $:
      \begin{align*}
            F'(n) &= \frac{\sum_{i=1}^{d_A} \lambda_i^n \ln \lambda_i }{\sum_{j=1}^{d_A} \lambda_j^n} 
            = \sum_i z_i(n) \ln \lambda_i \qquad \text{and} \\
            F''(n) &= \sum_i z_i(n) (\ln \lambda_i)^2 - \left( \sum_i z_i(n) \ln \lambda_i \right)^2 \\
            &= \mathrm{Var}_{z(n)}(\ln \lambda_i) \geq 0 ,
      \end{align*}
      which implies, that $F'$ is non-decreasing. Thus, we can follow w.l.o.g for $x\geq n$
      \begin{align*}
            F(x) - F(n) &= \int_n^x F'(s) ds \leq \int_n^x F'(n) ds = F'(n) (x - n),
      \end{align*}
      which is equivalent to saying that $F$ lays above all of its tangent, since it is convex. Since $F(1) = \ln \sum_i \lambda_i = \ln 1 = 0 $, it follows that
      \begin{equation*}
            F(1) = 0 \leq F(n) + F'(n)(1-n) .
      \end{equation*}
      We can employ this now to yield 
      \begin{equation*}
            \frac{d S(n)}{dn} = \frac{d}{dn} \frac{F(n)}{1-n} = \frac{(1-n) F'(n) + F(n)}{(1-n)^2} \leq 0 ,
      \end{equation*}  
      which proves, that $S_n$ is non-increasing for $n \neq 1$.
      Finally, employing l'Hôpital's rule yields for $n=1$ that
      \begin{equation*}
            \lim_{n\rightarrow1} = - \frac{1}{2} F''(1) \leq 0 .
      \end{equation*}
\end{claimproof}



\subsection{Computing von Neumann entropy with variational quantum states}
\label{subsec:von-neumann-entanglement-entropy-vmc}


\subsubsection{Direct computation using the SVD}

The entanglement entropy for the ground state can be computed directly as follows:
\begin{enumerate}
\item Obtain an approximation of the ground state
      $\ket{\psi_0}=\sum_\alpha c_\alpha\ket{\alpha}\in\mathcal H^{AB} = \mathcal H^{A} \otimes \mathcal H^{B} $ as some parametrized variational state $\psi_\theta(\alpha)$.
\item Use the variational state (e.g. $\psi_\theta(\alpha) = c_\alpha$) to evaluate the amplitudes on every allowed configuration $\alpha$ and normalize the resulting coefficient vector.
\item Reshape the coefficients into the matrix
      $\psi\in\mathbb C^{d_A\times d_B}$, respecting the ordering of the
      product basis.
\item Compute the SVD of $\psi$ to obtain the Schmidt values $s_i$, yielding the entropy $S=-\sum_i s_i^2\ln s_i^2$
\end{enumerate}
This approach is also straightforward with exact diagonalization since the ground state is represented by a vector consisting of the desired amplitudes. However, this quickly becomes unfeasible for ED and variational states as the Hilbert space grows exponentially with the system size.


\subsubsection{Constructing the reduced density matrix}

Given that $d_A\ll d_B$, one can estimate the small reduced density matrix $\rho_A$ and diagonalize it directly. For illustrative purposes, let subsystem $A$ describe a single photon mode and let subsystem $B$ describe a fermionic system.

Let $n=0,\dots,N_{\mathrm{ph}}$ denote the photon occupation and $\sigma$
a fermionic occupation configuration. Define the wavefunction amplitudes by
\begin{equation}
\psi(n,\sigma)\equiv\braket{n,\sigma}{\psi}.
\end{equation}
For a normalized state, the matrix elements of $\rho_A$ follow from \cref{eq:reduced_density_matrix_as_operator,eq:density_matrix_as_wavefunction_coefficients}
\begin{equation}
(\rho_A)_{nn'}
=\sum_\sigma\psi(n,\sigma)\psi^*(n',\sigma)
\end{equation}

Since variational states do not necessarily provide normalized amplitudes, we might have to normalize them in the above formula. Defining
\begin{equation}
Z\equiv\sum_{m,\sigma}|\psi(m,\sigma)|^2,
\end{equation}
we can write the normalized reduced density matrix more generally as
\begin{equation}
\boxed{\,
(\rho_A)_{nn'}
=\frac{\sum_\sigma\psi(n,\sigma)\psi^*(n',\sigma)}{Z}
\,}.
\end{equation}
For a normalized state, $Z=1$.

\subsubsection{A Monte Carlo estimator for the reduced density matrix}

For each fermionic configuration $\sigma$, define
\begin{equation}
u_n(\sigma)\equiv
\frac{\psi(n,\sigma)}
{\sqrt{\sum_{m=0}^{N_{\mathrm{ph}}}|\psi(m,\sigma)|^2}},
\qquad
\vec u(\sigma)=
\begin{pmatrix}
u_0(\sigma)\\
\vdots\\
u_{N_{\mathrm{ph}}}(\sigma)
\end{pmatrix}
\end{equation}
only if the denominator is non-zero. This is almost always the case when using Monte Carlo sampling since these configurations will never be visited, unless you start in one of them.
This vector is normalized for every $\sigma$ with non-zero marginal
probability:
\begin{equation}
\vec u(\sigma)^\dagger\vec u(\sigma)
=\sum_{n=0}^{N_{\mathrm{ph}}}|u_n(\sigma)|^2=1.
\end{equation}
The matrix elements of its outer product are
\begin{equation}
\bigl(\vec u(\sigma)\vec u(\sigma)^\dagger\bigr)_{nn'}
=
\frac{\psi(n,\sigma)\psi^*(n',\sigma)}
{\sum_{m=0}^{N_{\mathrm{ph}}}|\psi(m,\sigma)|^2}.
\end{equation}

Now consider the marginal probability distribution of the fermionic
configurations,
\begin{align}
p(\sigma) &=
\sum_{m=0}^{N_{\mathrm{ph}}} p(m, \sigma) \\
&\equiv \frac{\sum_{m=0}^{N_{\mathrm{ph}}}|\psi(m,\sigma)|^2}
{\sum_{m,\sigma'}|\psi(m,\sigma')|^2} .
\end{align}
Taking the expectation value with respect to this distribution yields
\begin{align}
\mathbb E_{\sigma\sim p(\sigma)}
\bigl[\vec u(\sigma)\vec u(\sigma)^\dagger\bigr]_{nn'}
&=\sum_\sigma
  \frac{\psi(n,\sigma)\psi^*(n',\sigma)}
       {\sum_m|\psi(m,\sigma)|^2}\,p(\sigma)\notag\\
&=\sum_\sigma
  \frac{\psi(n,\sigma)\psi^*(n',\sigma)}
       {\sum_m|\psi(m,\sigma)|^2}\,
  \frac{\sum_m|\psi(m,\sigma)|^2}
       {\sum_{m,\sigma'}|\psi(m,\sigma')|^2}\notag\\
&=\frac{\sum_\sigma\psi(n,\sigma)\psi^*(n',\sigma)}
       {\sum_{m,\sigma'}|\psi(m,\sigma')|^2}\notag\\
&=(\rho_A)_{nn'}.
\end{align}
Therefore,
\begin{equation}
\boxed{\,
\rho_A=
\mathbb E_{\sigma\sim p(\sigma)}
\bigl[\vec u(\sigma)\vec u(\sigma)^\dagger\bigr]
\,}.
\end{equation}
Moreover, since we are only interested in an expectation value over $\sigma$, we have
\begin{align}
\mathbb E_{(n,\sigma)\sim p(n, \sigma)}
\bigl[F(\sigma)\bigr]
&=\sum_{n,\sigma}
  p(n, \sigma) F(\sigma) \notag\\
&=\sum_{\sigma}
  p(\sigma) F(\sigma) \notag\\
&= \mathbb E_{\sigma\sim p(\sigma)}
\bigl[F(\sigma)\bigr] . \label{eq:expectation_value_equal_probability_marginal_probability}
\end{align}
Thus, we can sample either from $p(\sigma)$ or from $p(n,\sigma)$, depending on what is more convenient.

\subsubsection{Computation using variational Monte Carlo}

The resulting procedure is:
\begin{enumerate}
\item Sample a configuration $(\tilde n,\sigma)$ using the VMC sampler,
      i.e.\ from the joint distribution
      \begin{equation}
      p(n,\sigma)=\frac{|\psi(n,\sigma)|^2}{Z}.
      \end{equation}
      Discard $\tilde n$ and retain $\sigma$. Note, that $Z$ does not need to be known in VMC. The retained configurations follow the marginal distribution $p(\sigma)$ (see \cref{eq:expectation_value_equal_probability_marginal_probability}).
\item Evaluate all amplitudes $\psi(n,\sigma)$ for
      $n=0,\dots,N_{\mathrm{ph}}$ and construct $\vec u(\sigma)$.
\item Compute $\vec u(\sigma)\vec u(\sigma)^\dagger$.
      Repeat the preceding steps and average the resulting matrices:
      \begin{equation}
      \widehat\rho_A=
      \frac{1}{M}\sum_{k=1}^{M}
      \vec u(\sigma_k)\vec u(\sigma_k)^\dagger
      \approx\rho_A.
      \end{equation}
      The wide hat indicates that we only obtain an estimate of the true density matrix.
\item Diagonalize $\widehat\rho_A$ and use its eigenvalues
      $\widehat\lambda_i$ to estimate the von Neumann entropy:
      \begin{equation}
      \widehat S=-\sum_i\widehat\lambda_i\ln\widehat\lambda_i.
      \end{equation}
\end{enumerate}

\subsection{References}
I did not use any literature explicitly, except for AI helping me figure out what to do. The following paper employs a similar approach to the one I demonstrate here.
\begin{itemize}
      \item M.\ Nielsen \& I.\ Chuang, \emph{Quantum Computation and Quantum Information}
      (density operator, partial trace; Box 2.6).
      \item U.\ Schollw\"ock, \emph{The density-matrix renormalization group in the age of matrix
            product states}, beginning of Chapter 4 (MPS).
      \item Tang, Yifan, et al. \emph{Deep quantum Monte Carlo approach for polaritonic chemistry.} The Journal of Chemical Physics 163.3 (2025). (Appendix A.1 and A.3)
\end{itemize}

\end{document}
